\section{Quantum Entanglement}
When multiple qubits are combined, their joint state space is formed by the multiplicative tensor product of their individual vector spaces.
This specific mathematical structure naturally gives rise to collective system states that have absolutely no classical equivalent.

A multi-qubit state is mathematically defined as \textit{separable} (or a product state) if it can be cleanly factored into the tensor product of independent, individual qubit states.
For instance, the uniform superposition state $\ket{\psi} = \frac{1}{2}(\ket{00} + \ket{01} + \ket{10} + \ket{11})$ is highly parallel, but it is still fundamentally separable.
This is because it can be factored perfectly into $\left(\frac{1}{\sqrt{2}}(\ket{0} + \ket{1})\right) \otimes \left(\frac{1}{\sqrt{2}}(\ket{0} + \ket{1})\right)$, meaning manipulating or measuring one qubit does not mathematically force an outcome on the other.

However, the vast mathematical majority of states residing in a composite Hilbert space cannot be factored in this manner.
States that inherently resist this factoring are called \textit{entangled}.
Entanglement represents a profound form of non-local correlation, acting as the primary fuel for the massive speedups seen in advanced quantum algorithms.
In a purely entangled system, the quantum state of the collective whole is perfectly defined, yet the states of the individual constituent subsystems remain completely undefined.
This non-classical phenomenon was famously, and somewhat skeptically, described by Albert Einstein, Boris Podolsky, and Nathan Rosen (EPR) as "spooky action at a distance" \cite{einstein1935}.

\subsection{The Bell States}
The most iconic examples of maximally entangled two-qubit systems are the four Bell states, often referred to in the literature as EPR pairs \cite{nielsen2010quantum}.
They form a complete, orthonormal basis for the entire two-qubit state space and are mathematically defined as follows:
\begin{align}
    \ket{\Phi^+} &= \frac{1}{\sqrt{2}}(\ket{00} + \ket{11}) \\
    \ket{\Phi^-} &= \frac{1}{\sqrt{2}}(\ket{00} - \ket{11}) \\
    \ket{\Psi^+} &= \frac{1}{\sqrt{2}}(\ket{01} + \ket{10}) \\
    \ket{\Psi^-} &= \frac{1}{\sqrt{2}}(\ket{01} - \ket{10})
\end{align}

Consider the specific entangled state $\ket{\Phi^+}$.
Imagine two experimental parties, traditionally named Alice and Bob, each holding exactly one isolated qubit of this pair.
If Alice measures her isolated qubit, she inherently faces a 50\% probability of observing a $0$ and a 50\% probability of observing a $1$.
However, the very moment she observes a $0$, the global wave function of the entire bipartite system instantaneously collapses from the superposition into the definite state $\ket{00}$.

Consequently, if Bob subsequently measures his corresponding qubit, he will observe a $0$ with absolute 100\% certainty.
This correlation holds true regardless of whether Bob is sitting in the same laboratory as Alice or is separated by light-years across the galaxy.

This perfect, instantaneous correlation statistically violates the mathematical limits established by Bell's inequalities.
This violation is deeply significant because it definitively proves that quantum mechanics cannot be explained by "local hidden variable" theories \cite{bell1964}.
In other words, the measurement outcomes are not secretly pre-programmed into the particles before they separate; the universe is fundamentally, irreducibly probabilistic at the quantum level.

Because an entangled system cannot be mathematically decoupled into independent individual states, the single-qubit Bloch sphere representation fails entirely when attempting to visualize entanglement.
The fundamental physical paradox here is one of information distribution.
While we can possess maximal knowledge about the joint system as a whole (a pure state), we inherently possess absolutely zero specific knowledge about its individual constituent parts prior to measurement.